% H*(MO∧MO; F2) as a module over the Steenrod algebra at p = 2: four towers of
% classes joined by h0 (vertical), Sq1 (short arc) and Sq2 (long arc), with
% dashed lines for the sum decompositions.
% Rows are counted from the top (row 1) to the bottom (row 10); column c
% occupies rows top_c..10.
\documentclass[tikz,border=2pt]{standalone}
\usepackage{dzg-tikz}

\begin{document}
\begin{tikzpicture}[x=3.2cm, y=-0.5cm]
  \foreach \c/\top in {1/1, 2/2, 3/4, 4/5} {
    \foreach \r in {\top,...,10}
      \node[dual vertex] (d-\r-\c) at (\c,\r) {};
    % Operations from row r reach rows r+1 (h0), r+2 (Sq1), r+3 (Sq2).
    \foreach \r in {\top,...,9} {
      \pgfmathtruncatemacro{\h}{\r+1}
      \pgfmathtruncatemacro{\one}{\r+2}
      \pgfmathtruncatemacro{\two}{\r+3}
      \draw (d-\r-\c) -- (d-\h-\c);
      \ifnum\one<11 \draw (d-\r-\c) to[bend left=35] (d-\one-\c); \fi
      \ifnum\two<11 \draw (d-\r-\c) to[bend left=55] (d-\two-\c); \fi
    }
  }
  \node[left, font=\small] at (d-2-1) {$w_2U\otimes U$};
  \node[left, font=\small] at (d-6-1) {$U\otimes w_1U$};

  % Sum decompositions.
  \draw[dashed] (d-10-2) -- (d-6-1);
  \draw[dashed] (d-10-4) -- (d-8-3);

  % Column labels.
  \foreach \c/\lbl in {1/{U\otimes U}, 2/{U\otimes w_1^2U + w_2U\otimes U},
      3/{w_2^2U\otimes U}, 4/{w_2^2U\otimes w_1^2U + w_2^3U\otimes U}}
    \node[below=1.5em, font=\small] at (d-10-\c) {$\lbl$};
\end{tikzpicture}
\end{document}
